\section{Guerrilla Units}\label{sec:7.0}

\ruleslabel{General Rule}

\begin{generalrule}
Throughout the game, the Yugoslav Player controls three types of units: guerrillas, Soviets, and pro-Yugoslav Bulgarian units. Soviet and Bulgarian units are available only as reinforcements on Game-Turn 15. Thus, from Game-Turns 1 through 14, the Yugoslav Player controls only guerrilla units. Guerrilla units are of two types: Partisan (those guerrillas controlled by Tito) and pro-Yugoslav Chetnik (Serbian royalists). It is possible that Chetnik units will be pro-Yugoslav, pro-Axis, or neutral (see 7.2). Guerrilla units come in three sizes: groups, brigades, and Partisan divisions. However, there are specific restrictions concerning the time in which brigades and divisions may be first created by the Yugoslav Player. At the beginning of the game, all guerrilla units operate in groups.
\end{generalrule}

\ruleslabel{Cases}

\subsection{Partisans}\label{sec:7.1}

Partisan units (including the Tito counter) are always under the control of the Yugoslav Player.

\subsection{Chetniks}\label{sec:7.2}

\case{7.21} During each Chetnik Collaboration Phase between Game-Turns 2 (inclusive) and the Game-Turn in which Allied Support Withdrawal takes place (exclusive; see 7.26), the Yugoslav Player must roll a single die for each box, triangle, or circle currently occupied by at least one Chetnik unit (including the Neutral Chetnik boxes):

\textbf{A.} If the die roll is 2 or less, the Chetnik units become (or remain) pro-Axis (see 7.24);

\textbf{B.} If the die roll is 3 or 4, the Chetnik units become (or remain) neutral (see 7.24);

\textbf{C.} If the die roll is 5 or more, the Chetnik units become (or remain) pro-Yugoslav (see 7.24).

\case{7.22} During each Chetnik Collaboration Phase between the Game-Turn in which Allied Support Withdrawal (inclusive) takes place and the end of the game (inclusive), the Yugoslav Player must roll a single die for each box, triangle, or circle currently occupied by at least one Chetnik unit:

\textbf{A.} If the die roll is 2 or less, the Chetnik units become (or remain) pro-Axis (see 7.24);

\textbf{B.} If the die roll is 3 through 6, the Chetnik units disband (they are immediately removed from the map);

\textbf{C.} If the die roll is 7 or more, the Chetniks become Partisan. They are immediately replaced by an equal number and type of Partisan units, which are placed in the Yugoslav/Partisan Mountain of the occupied Zone. If an equal number and type of Partisan units are not available for use, the Chetniks are considered disbanded (see 7.22B).

\case{7.23} The following modifiers are applied to all Chetnik Collaboration die rolls (all modifiers are cumulative):

\textbf{A.} +1: If the Chetnik stack is currently pro-Yugoslav;

\textbf{B.} +1: If an Italian Surrender has taken place (see 10.3);

\textbf{C.} +1: In all Game-Turns after and including the Game-Turn in which Allied Support Withdrawal takes place (see 7.26);

\textbf{D.} $-1$: If the Chetnik stack is currently pro-Axis;

\textbf{E.} $-1$: If Tito has been eliminated (see 9.4).

\case{7.24} If Chetnik units become pro-Axis, they pass to the control of the Axis Player at the end of the Phase by being transferred to the Axis Mountain triangle of the occupied Zone’s Zone Display. If Chetnik units become pro-Yugoslav, they pass to the control of the Yugoslav Player at the end of the Phase by being transferred to the Yugoslav/Chetnik Mountain triangle of the occupied Zone’s Zone Display. If Chetnik units become neutral, they are placed in the occupied Zone’s Neutral Chetnik box at the end of the Phase. Here they remain, uncontrolled by either Player, until future die rolls make them pro-Axis or pro-Yugoslav.

\case{7.25} If Chetnik units remain pro-Axis, pro-Yugoslav, or neutral, they are not moved from their current location.

\case{7.26} At the beginning of the Chetnik Collaboration Phases of Game-Turns 9 through 13, the Yugoslav Player must determine if Allied Support Withdrawal has taken place. Allied Support Withdrawal occurs only once per game, and after its initiation, it is in effect for the remainder of the game (i.e., it is unnecessary to calculate its occurrence again in succeeding Game-Turns).

\textbf{A.} On Game-Turn 9, Allied Support Withdrawal occurs if less than 20\% of the Chetnik units currently on the map are pro-Yugoslav;

\textbf{B.} On Game-Turn 10, Allied Support Withdrawal occurs if less than 40\% of the Chetnik units on the map are pro-Yugoslav;

\textbf{C.} On Game-Turn 11, Allied Support Withdrawal occurs if less than 60\% of the Chetnik units on the map are pro-Yugoslav;

\textbf{D.} On Game-Turn 12, Allied Support Withdrawal occurs if less than 80\% of the Chetnik units on the map are pro-Yugoslav;

\textbf{E.} On Game-Turn 13, Allied Support Withdrawal occurs automatically if it has not occurred already.

\case{7.27} It is not necessary to perform Chetnik Collaboration die rolls when there are no Chetnik units on the map.

\subsection{Guerrilla Movement And Combat Restrictions}\label{sec:7.3}

\case{7.31} Partisan and pro-Yugoslav Chetnik units may never occupy the same box, triangle, or circle. They may never participate in an attack together and may never be attacked as a combined force.

\case{7.32} Pro-Axis Chetnik units may stack freely with other Axis units. When stacked with Axis units that are performing combat, they must participate. Exception: See 7.33.

\case{7.33} Pro-Yugoslav Chetnik and pro-Axis Chetnik units may never perform combat against one another. Pro-Axis Chetnik units are ignored when stacked with Axis units that are performing combat against pro-Yugoslav Chetniks. Pro-Yugoslav Chetniks may attack Axis units stacked with pro-Axis Chetniks (although the pro-Axis Chetniks are ignored in the attack).

\subsection{Guerrilla Reinforcements}\label{sec:7.4}

During the course of the game, new guerrilla units may be created through pre-determined reinforcements or variable reinforcements.

\case{7.41} The Yugoslav Player receives predetermined Partisan reinforcements (including Tito) at the beginning of the Yugoslav Movement Phase of Game-Turn 2. These units are placed according to the instructions of Case 13.92.

\case{7.42} During all Guerrilla Reinforcement Phases, the Yugoslav Player must determine if variable reinforcements are available. In each of the 12 Zones, the Yugoslav Player must determine if reinforcements are available through recruitment (see 7.43), due to Tito (see 7.46), or through a Guerrilla Uprising (see 7.47). Reinforcements are expressed in groups (see 7.6) of Partisans and/or Chetniks.

\case{7.43} There are two means by which the Yugoslav Player can obtain reinforcements through recruitment: by occupying an Objective Display (see 7.44) or by occupying a Mountain triangle (see 7.45). Note that recruitment may be affected by Tito’s removal from the map or his elimination (see 9.32) or by weather (see 11.0).

\case{7.44} Each Objective Display currently occupied by at least one Yugoslav unit (including Soviets and pro-Yugoslav Bulgarians) is eligible for recruitment during the Recruitment Segment. In order to determine the number of groups created due to a Yugoslav-occupied Display, the Yugoslav Player rolls a single die and multiplies this die roll by the Display’s Recruitment Value (round down if a fraction is created). The result is the number of reinforcing groups created. These reinforcements are Partisan if the Display is currently occupied by a Partisan (or Soviet or pro-Yugoslav Bulgarian) stack, or Chetnik if the Display is currently occupied by a Chetnik stack. The number of reinforcements created on a Display is subject to the following modifications:

\textbf{A.} No more reinforcing groups may be created than the number of Yugoslav Strength Points currently occupying the Display;

\textbf{B.} If a drought is in effect (see 11.2), the number of groups created in market town Objective Displays is halved (round fractions down);

\textbf{C.} If a Winter Game-Turn is in effect (see 11.1), the number of groups created in village Objective Displays is halved (round fractions down);

\textbf{D.} If a Zone’s Alignment Value is Partisan (see 7.5), the number of reinforcing Chetnik groups created by recruitment in all Objective Displays within that Zone is halved (round fractions down). If a Zone’s Alignment Value is Chetnik (see 7.5), the number of reinforcing Partisan groups created by recruitment in all Objective Displays within that Zone is halved (round fractions down). If a Zone’s Alignment Value is neutral, recruitment in Objective Displays within that Zone is unaffected.

\case{7.45} Each Mountain triangle currently occupied by at least one Yugoslav unit is eligible for recruitment in the Recruitment Segment. In order to determine the number of groups created due to this cause, the Yugoslav Player should total the number of Yugoslav Strength Points occupying a given Mountain triangle and divide by four (round fractions down). The result is the number of reinforcing groups created within that Mountain. Reinforcing groups are Partisan if the Mountain triangle is occupied by Partisan (or Soviet or pro-Yugoslav Bulgarian) units. Reinforcing groups are Chetnik if the Mountain triangle is occupied by Chetnik units. The number of reinforcements created within a Mountain is subject to the following modifications:

\textbf{A.} No more groups may be created than the number of Yugoslav units currently occupying a given Mountain triangle.

\textbf{B.} If it is a Winter Game-Turn (see 11.1), the number of groups created in each Mountain triangle is halved (round fractions down).

\textbf{C.} If a Zone’s Alignment Value is Partisan (see 7.5), the number of reinforcing Chetnik groups created by recruitment in the Mountain triangle within that Zone is halved (round fractions down). If a Zone’s Alignment Value is Chetnik (see 7.5), the number of reinforcing Partisan groups created by recruitment in the Mountain triangle within that Zone is halved (round fractions down). If a Zone’s Alignment Value is neutral, recruitment in Mountains within that Zone is unaffected. Note: The presence of Axis units in a corresponding triangle has no effect on recruitment.

\case{7.46} During the Tito Segment, the Yugoslav Player determines if Tito (in
any state) is present anywhere within this Zone, including the Hide-away circle.
If Tito is present, the Yugoslav Player receives a number of Partisan groups as
reinforcements according to the following schedule:

\textbf{A.} If the Zone is Croatia, 3 groups are created.

\textbf{B.} If the Zone is Baranya, Bosnia, Carinthia, Dalmatia, or Slovenia, 2 groups are created.

\textbf{C.} If the Zone is not listed above, only 1 group is created.

\case{7.47} During the Uprising Segment, the Yugoslav Player determines if any Zones are eligible for a Guerrilla Uprising (see 7.5). If a Guerrilla Uprising is successful, a number of Partisan and/or Chetnik Strength Points will be created.

\case{7.48} Variable guerrilla reinforcements are deployed as follows:

\textbf{A.} If they are created by recruitment in an Objective Display, they are immediately placed in the Yugoslav box of that Display.

\textbf{B.} If they are created by recruitment in a Mountain triangle, they are immediately placed directly in that Mountain.

\textbf{C.} If they are created by Tito’s presence within a Zone, they are immediately placed directly on top of the Tito counter.

\textbf{D.} If they are created by a Guerrilla Uprising, they are immediately
placed in the appropriate Hide-away circle of the Zone in which the uprising
occurred.

\case{7.49} Anytime the number of reinforcing guerrilla groups created by recruitment is halved (do not count Recruitment Values of \textonehalf) more than once in a given Objective Display or Mountain triangle, no reinforcements are created during the current Recruitment Segment for the Display or triangle.

Example of Variable Reinforcements:

During the Guerrilla Reinforcement Phase, the Yugoslav Player is calculating reinforcements in Slovenia. One Partisan Division (Strength of 12) plus Tito occupies the Partisan Mountain triangle. Nine Chetnik groups (each with a Strength of 1) occupy the Chetnik Mountain triangle. In addition, one Partisan group occupies Ljubljana. During the Recruitment Segment, the Yugoslav Player rolls a single die for Ljubljana, obtaining a 6. Multiplied by Ljubljana’s Recruitment Value (2), a 12 results. However, only one Partisan group is created, because only a single Strength Point occupies this Display (see 7.44A). This group is immediately deployed in Ljubljana (if available). For the Mountains, only a single Partisan group is created due to the Partisan division (12 divided by 4 equals 3; since there is only one Partisan unit in the triangle, only a single group is created — see 7.45). Similarly, one Chetnik group is created by the Chetnik presence in the Mountain triangle (9 divided by 4 equals 2\textonequarter, dropping fractions leaves 2; 2 is halved to 1 because Slovenia is aligned to the Partisans — see 7.45C). Both of these reinforcing groups are deployed in their respective Mountain triangles. During the Tito Segment, the Yugoslav Player receives 2 Partisan groups due to Tito’s presence in the Zone (see 7.46), both of which are placed on top of Tito in the Partisan Mountain triangle. During the Uprising Segment, the Yugoslav Player determines if an uprising takes place in Slovenia. However, the Axis garrison is sufficient to prevent its occurrence (see 7.5).

\subsection{Guerrilla Uprisings}\label{sec:7.5}

A Guerrilla Uprising may take place within a Zone if the Axis garrison falls below the minimum strength required to keep the Zone pacified. On the map, the name of each Zone is followed by three Values. The first is the Garrison Value, the Second (a letter) is the Alignment Value, and the third is the Uprising Modifier.

\case{7.51} If, at the beginning of the Uprising Segment of all Game-Turns between 3 and 17 (inclusive), the Axis Player currently has fewer divisions (or their equivalents; see 6.22) within a Zone than that Zone’s Garrison Value, a Guerrilla Uprising automatically takes place.

\case{7.52} In each Zone in which a Guerrilla Uprising takes place, the Yugoslav Player immediately rolls a single die. The Zone’s Uprising Modifier is immediately subtracted from this die roll. A modified die roll of 1 or more indicates a successful Guerrilla Uprising (see 7.53). A modified die roll of less than 1 indicates an unsuccessful Guerrilla Uprising. If a Guerrilla Uprising is unsuccessful, nothing further is done.

\case{7.53} If a Guerrilla Uprising is successful, the following procedure is performed:

\textbf{A.} The Axis Player determines the difference between the Garrison Value of the Zone and the actual number of divisions currently occupying this Zone (drop fractions of divisions);

\textbf{B.} The figure determined in Step A is multiplied by the Yugoslav Player’s original modified die roll which created the successful uprising (see 7.52). This product is the number of reinforcing groups created as a result of the uprising;

\textbf{C.} The Yugoslav Player rolls a single die and compares this roll with the Zone’s Alignment Value (either P - Partisan; or C - Chetnik; or N - Neutral). If an even number is rolled, all of the reinforcing groups created as a result of the uprising are of the same type as the Zone’s alignment. If the Zone is neutral, no die roll is necessary and half the groups are Partisan (round fractions up) and the other half are Chetnik. If an odd number is rolled, half the reinforcing groups are Partisan (round fractions up) and the other half are Chetnik, regardless of the Alignment Value of the Zone.

\case{7.54} Axis units on an Anti-Guerrilla Operation within a Zone (see 8.4) may be fully counted as part of that Zone’s garrison, even though they are not currently deployed on the map.

\case{7.55} No Guerrilla Uprisings may take place if there are no guerrilla units on the map.

\case{7.56} Reinforcements created by a Guerrilla Uprising are placed in the appropriate Hide-away circle of the Zone in which the uprising occurred.

Example of Guerrilla Uprising:

In Bosnia, the Axis Player has nine Croat brigades. This is equal to 4½ divisions — below Bosnia’s Garrison Value of 6. Thus, a Guerrilla Uprising takes place. The Yugoslav Player rolls a single die and obtains a 2. Bosnia’s Uprising Modifier is 1, so the modified die roll is 1 — a successful uprising. The Axis Player is two divisions short of the Garrison Value, so a 2 multiplied by 1 yields 2. The Yugoslav Player again rolls the die, obtaining a 4 (even). Thus, the two reinforcing guerrilla groups created in Bosnia are both Partisan (Bosnia’s Alignment Value is P). These groups are placed in the Partisan Hide-away in Bosnia.

\subsection{Maximum Size Of Guerrilla Units}\label{sec:7.6}

Guerrilla units come in three sizes: groups, brigades, and divisions (Partisan only). Brigades and divisions are created from smaller guerrilla units. When creating brigades and divisions, any unused brigade or division counter may be employed.

\case{7.61} At the beginning of the game and in all succeeding Game-Turns until brigade-strength is achieved (see 7.62), the maximum size of any guerrilla unit on the map is the group.

\case{7.62} At the end of each Guerrilla Status Phase, the Yugoslav Player should total the number of Partisan and the number of Chetnik groups currently on the map:

\textbf{A.} If there are 30 Partisan groups on the map at this time, the Yugoslav Player has achieved brigade-strength for the Partisans;

\textbf{B.} If there are 25 Chetnik groups of any allegiance on the map at this time, the Yugoslav Player has achieved brigade-strength for the Chetniks.

Note: There are 30 and 25 Partisan and Chetnik groups provided in the counter mix, respectively. Thus, when all of the groups of a particular guerrilla type are on the map, brigade-strength has been achieved for that type.

\case{7.63} If the Partisans (or Chetniks) have achieved brigade-strength, the Yugoslav Player may build Partisan (or Chetnik) groups into brigades, subject to the restrictions of Cases 7.64 and 7.67:

\textbf{A.} Immediately (and in every succeeding Guerrilla Status Phase), the Yugoslav Player may create brigades from groups of the same guerrilla type that currently occupy the same box, triangle or circle. 3 groups are the equivalent of 1 brigade. In each Guerrilla Status Phase, the Yugoslav Player may perform as many of these substitutions as he

desires. Substituted units are removed from the map and replaced by an unused brigade counter. No larger unit may be created in a Hide-away.

\textbf{B.} In ensuing Guerrilla Reinforcement Phases, the Yugoslav Player may bring reinforcing Partisan or Chetnik groups onto the map in brigade size. 3 groups are the equivalent of 1 brigade. Example: If in Zagreb the Yugoslav Player were due 7 Partisan groups as reinforcements, 2 brigades and 1 group, or 1 brigade and 4 groups, or simply 7 groups could be deployed.

\case{7.64} Only Chetnik units that are controlled by the Yugoslav Player may be built into brigade-size. Partisan and Chetnik units may never combine to form any higher level of organization.

\case{7.65} At the end of each Guerrilla Status Phase after brigade-strength has been achieved for the Partisans, the Yugoslav Player should total the number of Partisan brigades currently on the map. If this number is 25 or more, division-strength has been achieved by the Partisans.

\case{7.66} If the Partisans have achieved division strength, the Yugoslav Player may create Partisan divisions according to the restrictions of Case 7.67. Chetniks may never form into divisions.

\textbf{A.} Immediately (and in every succeeding Guerrilla Status Phase), the Yugoslav Player may create divisions from groups and/or brigades that occupy the same box, triangle, or circle. 2 brigades or 6 groups are the equivalent of 1 division. In each Guerrilla Status Phase, the Yugoslav Player may perform as many of these substitutions as he desires. Substituted units are removed from the map and replaced by an unused division counter. No larger unit may be created in a Hide-away.

\textbf{B.} In ensuing Guerrilla Reinforcement Phases, the Yugoslav Player may bring Partisan groups onto the map in division-size. 6 groups or 2 brigades are the equivalent of 1 division.

\case{7.67} The first Partisan brigade and division that is placed on the map must be created in the box, triangle, or circle that Tito currently occupies. After this requirement is fulfilled, Partisan brigades and divisions may be created for the remainder of the game without further restriction. If Tito is removed from the map or eliminated (see 9.4) at the time brigade or division-strength is achieved, brigades or divisions may not be created until Tito’s redeployment on the map.

\case{7.68} The number of Yugoslav groups, brigades, and divisions provided in Tito is an intended limitation. If no more counters of a given type are available for use, then no more units of this type may be created by any means of reinforcement. However, all guerrilla units that are eliminated immediately become available for use again upon their elimination. Note that all Yugoslav groups and divisions are back-printed with brigade-size units. Similarly, Italian units are back-printed with Partisan divisions; as Italian units are eliminated from play or are withdrawn from the map, the Yugoslav Player may employ the reverse side of these counters as divisions become available for employment.

\case{7.69} Once the Yugoslav Player is permitted to build brigades or divisions,
he may continue to do so for the remainder of the game, even if conditions for
their creation may have fallen below their original ‘‘triggering’’ point.
However, once a Yugoslav brigade or division is created, it may never be broken
down again for any reason (including combat loss) for the remainder of the game.
Note: The achievement of brigade-strength for Partisans or Chetniks may trigger
a number of key events (see 6.64 and 9.4).
